% Copyright (c) 2021-2027 Patrick Hall

\documentclass[fleqn]{article}
\renewcommand\refname{}
\title{Responsible Machine Learning\\\Large{Assignment 7}\\\Large{40 points}}

\usepackage{graphicx}
\usepackage{fullpage}
\usepackage{pdfpages}
\usepackage{amsmath}
\usepackage{amssymb}
\usepackage{mathtools}
\usepackage{MnSymbol}
\usepackage{enumerate}
\usepackage{setspace}
\usepackage[colorlinks, breaklinks=true]{hyperref}
\usepackage{float}
\usepackage{caption}
\usepackage{subcaption}
\usepackage{multicol}
\usepackage{color}
\usepackage{listings}
\usepackage{csvsimple}
\usepackage{algorithm}
\usepackage{algorithmic}
\usepackage{verbatim}
\usepackage{mdframed}
\usepackage{changepage}
\usepackage[top=1in, bottom=1in, left=1in, right=1in]{geometry}

\begin{document}

\maketitle

\noindent In Assignment 7, work alone or with a group to create a model card for your best model from previous assignments. You may use AI tools, but you are responsible for validating and understanding your work.

\section{Recommended References and Examples.}

Model cards document a model's intended uses, development, evaluation, and limitations. They help data scientists communicate about models so that others can maintain them and respond to incidents. Model documentation can also support executive review and may be required by regulation.\\

\noindent\textbf{Recommended references}:
\begin{itemize}
\item \href{https://arxiv.org/pdf/1810.03993.pdf}{\textit{Model Cards for Model Reporting}}
\item \href{https://deepmind.google/models/model-cards/}{Model Cards (Google DeepMind)}
\end{itemize}

\noindent\textbf{Example}: \url{https://github.com/jphall663/GWU_DNSC_6301_project}. This example is incomplete but may help you get started.

\section{Instructions and Rubric.}

Keep the model card concise. You may use bullet points, charts, and tables; number and describe each chart or table. The model card should focus on your best-performing model from previous assignments. Discuss additional models only as points of comparison, and explain why you selected the primary model.\\

\noindent\textbf{Rubric}:

\begin{itemize}
\item Structure (12 pts.)
\begin{itemize}
\item Name and contact information for yourself or all group members. You may anonymize this information for the assignment. (2 pts.)
\item Clearly delineated sections for:
\begin{itemize}
\item Intended use (1 pt.)
\item Training data (1 pt.)
\item Evaluation data (1 pt.)
\item Model details (1 pt.)
\item Quantitative analysis (1 pt.)
\item Ethical considerations (1 pt.)
\end{itemize}
\item Correct and informative external links, such as links to training and evaluation data (1 pt.)
\item Correct use and introduction of abbreviations (1 pt.)
\item Meaningful and informative charts or tables (1 pt.)
\item Proper descriptions of charts or tables (1 pt.)
\end{itemize}

\item Content (26 pts.)
\begin{itemize}
\item Intended use (4 pts.)
\begin{itemize}
\item Describe the business value of your best model.
\item Describe how the model is designed to be used and identify its intended users.
\item State whether the model can be used for additional purposes.
\end{itemize}

\item Training data (4 pts.)
\begin{itemize}
\item State the source of the training data and how it was divided into training and validation data.
\item State the number of rows in the training and validation data.
\item Define all training data columns and any engineered columns.
\end{itemize}

\item Evaluation data (4 pts.)
\begin{itemize}
\item State the source and number of rows of the evaluation (or test) data.
\item Describe any differences in columns between the training and evaluation data.
\end{itemize}

\item Model details (4 pts.)
\begin{itemize}
\item State the input and target columns used by your best model.
\item State the model type, software and software version, and the model's hyperparameters or other settings.
\end{itemize}

\item Quantitative analysis (6 pts.)
\begin{itemize}
\item State the metrics used to evaluate your best model (1 pt.).
\item Report these metrics for the training, validation, and evaluation (or test) data (1 pt.).
\item Submit your predictions to the class evaluation server at a TBD URL and report your model's leaderboard rank (1 pt.).
\item Include at least one plot or table from each weekly assignment, for a total of at least six (1 pt.).
\item Include global variable importance and partial dependence for your best model (1 pt.).
\item Briefly compare additional models as alternatives and explain why you selected your best model (1 pt.).
\end{itemize}

\item Ethical considerations (4 pts.)
\begin{itemize}
\item Describe potential negative impacts and uncertainties related to using your best model.
\item Consider math or software problems and real-world risks, including who could be affected, what could happen, and when and how it might occur.
\item Describe any unexpected results encountered during training.
\end{itemize}
\end{itemize}

\item Grammar and spelling (2 pts.)
\begin{itemize}
\item Correct grammar and punctuation (1 pt.)
\item Correct spelling (1 pt.)
\end{itemize}
\end{itemize}

\section{Submit Code and Model Card.}

Put your model card in the \texttt{README.md} file of a public GitHub repository, or a private repository shared with the instructor (\texttt{jphall663}). Submit your model card to the class evaluation server at a TBD URL.\\

\noindent\textbf{Your deliverables are due: TBD.}

\end{document}
